\section{Connectivness of Graphs}

\begin{definition}

                Let $G = (V, E)$ be a graph, $u, v \in V, u \neq v$. $u$ and $v$ are called \textbf{connected} if there is
                a $u-v$-path in $G$.
            
\end{definition}

\begin{definition}

                A graph $G = (V, E)$ is called \textbf{connected} if $\forall u, v \in V, u \neq v$ a $u-v$-path in $G$
                exists.
            
\end{definition}

\begin{definition}

                Let $G = (V, E)$ be a graph, $u, v \in V, u \neq v$. $u$ and $v$ are called \textbf{$k$-connected} if $|V|
                \ge k+1$ and
                $$\forall X \subseteq V \setminus \{u, v\} \text{ with } |X| \lt k, u \text{ and } v \text{ are
                connected in } G[V \setminus X]$$
            
\end{definition}

\begin{definition}

                A graph $G = (V, E)$ is called \textbf{$k$-connected} if $|V| \gt k$ and
                $$\forall X \subseteq V \text{ with } |X| \lt k, G[V \setminus X] \text{ is connected}$$
            
\end{definition}

\begin{definition}

                Let $G = (V, E)$ be a graph, $u, v \in V$ and $X \subseteq V \setminus \{u, v\}$. $X$ is called a
                \textbf{$u-v$-separator} if $u$ and $v$ are in different connected components of $G[V \setminus X]$.
            
\end{definition}

            More generally one can say that if a Graph $G$ contians a vertex $v$ stuch that $deg(v) \lt k$, then $G$ can
            not be $k$-connected.

            \begin{theorem}

\textbf{Menger's Theorem (Vertex Version)}:\\

                Let $G = (V, E)$ be a graph and $u, v \in V$. The minimum size of a $u-v$-separator is equal to the
                maximum number of internally vertex-disjoint $u-v$-paths.
                \\
\\

                A graph $G$ is $k$-connected iff. $\forall u, v \in V, u \neq v$ there are $k$ internally
                vertex-disjoint $u-v$-paths.
            
\end{theorem}

\begin{definition}

                A graph $G = (V, E)$ is called \textbf{$k$-edge-connected} if
                $$\forall X \subseteq E \text{ with } |X| \lt k, (V, E \setminus X) \text{ is connected}$$
            
\end{definition}

\begin{theorem}

\textbf{Menger's Theorem (Edge Version)}:\\

                A graph $G$ is $k$-edge-connected iff. $\forall u, v \in V, u \neq v$ there are $k$ edge-disjoint
                $u-v$-paths.
            
\end{theorem}

\begin{theorem}

                It always holds that:
                $$\text{Vertex Connectivity} \le \text{Edge Connectivity} \le \text{Minimum Degree } \delta(G)$$
            
\end{theorem}

            Lets look at the special case where $k=1$

            \begin{definition}

                Let $G = (V, E)$ be a connected graph.
                \begin{itemize}
    \item A vertex $v \in V$ is called a \textbf{cut vertex} (Artikulationsknoten) iff. $G[V \setminus
                        \{v\}]$ is disconnected.
    \item An edge $e \in E$ is called a \textbf{bridge} (Brücke) iff. $(V, E \setminus \{e\})$ is
                        disconnected.
\end{itemize}

\end{definition}

\begin{lemma}

                If $e = \{x, y\} \in E$ is a bridge, then $\text{deg}(x) = 1$ or $x$ is a cut vertex (analogously for
                $y$).
            
\end{lemma}

\begin{definition}

\textbf{Blocks}:\\

                We define an equivalence relation $\sim$ on $E$ by:
                $$e \sim f \iff e=f \lor \exists \text{ cycle containing } e \text{ and } f$$
                The equivalence classes are called \textbf{blocks}.
            
\end{definition}

            Note that two blocks can share at most one vertex, more over this vertex has to be a cut vertex.

            \begin{definition}

\textbf{Block Graph}:\\

                The block graph of a connected graph $G$ is a bipartite graph $T = (A \cup B, E_T)$ where:
                \begin{itemize}
    \item $A = \{\text{cut vertices of } G\}$
    \item $B = \{\text{blocks of } G\}$
    \item Edge $\{a, b\} \in E_T \iff a$ is incident to an edge in block $b$.
\end{itemize}

\end{definition}

\begin{theorem}

                If $G$ is connected, the block graph of $G$ is a tree.
            
\end{theorem}

\begin{algorithm}

\textbf{Depth-First Search (DFS) \& Low Values}:\\

                For efficient computation of cut vertices and bridges in time $O(|V|+|E|)$.
                \\

                $low[v] :=$ the smallest DFS number reachable from $v$ by a directed path consisting of tree edges and
                at most one back edge.
                \\

                It holds that:
                $$low[v] = \min \left\{ df s[v], \min_{(v,w) \text{ back edge}} df s[w], \min_{(v,w) \text{ tree edge}}
                low[w] \right\}$$
                A vertex $v$ is a cut vertex iff:
                \begin{enumerate}
    \item $v \neq \text{root}$ and there exists a child $u$ with $low[u] \ge df s[v]$.
    \item $v = \text{root}$ and has $\ge 2$ children in the DFS tree.
\end{enumerate}

                A tree edge $(v,w)$ is a bridge iff $low[w] \gt df s[v]$.
            
\end{algorithm}

